\clearpage\thispagestyle{empty}
\pdfbookmark[0]{Quellen zu den Beobachtungen der Parthenope}{parthenope719-sources}
\begin{center}\textbf{Quellen zu den Beobachtungen der Parthenope}\end{center}
\noindent Diese KI-Übertragung ins heutige Deutsch folgt dem Auszug auf S.~478 der \href{https://gdz.sub.uni-goettingen.de/id/PPN236007467}{\textit{Werke}, Band~VI} (Göttingen, 1874). Heinrich Christian Schumacher berichtet über die Beobachtungen der Parthenope, die Carl Friedrich Gauß am 31. Mai und 2. Juni 1850 in Göttingen machte.
\par\smallskip\noindent Die Beobachtungen stehen unter \textit{Beobachtungen und Elemente der Parthenope} in den \textit{Astronomischen Nachrichten}, Band~XXX, Nr.~719, \href{https://api.digitale-sammlungen.de/iiif/image/v2/bsb10049455_00191/full/1800,/0/default.jpg}{Sp.~374}. Die \href{https://api.digitale-sammlungen.de/iiif/image/v2/bsb10049455_00192/full/1800,/0/default.jpg}{Schlusszeile auf Sp.~376} nennt den 20. Juni 1850. Die Überschrift gehört zu einem längeren Zeitschriftenbericht. Die Werkausgabe gibt daraus nur die Göttinger Beobachtungen wieder, auf die in der Zeitschrift die von Georg Rümker berechneten Elemente folgen. KI-Übertragung und KI-Satz.
\clearpage
